%! Logarithmic Expressions — Unit 2, Lesson 2.9 — Lesson Plan
\documentclass[10pt]{article}
\usepackage{precalculus-article}
\usepackage{precalculus-boxes}
\usepackage{graphicx}
\graphicspath{{images/}}
\newcommand{\TallMath}[1]{$\displaystyle #1\rule[-1.4em]{0pt}{3.2em}$}
\newcommand{\CourseName}{Precalculus}
\newcommand{\SchoolYear}{2026--2027}
\newcommand{\MeetingLength}{55 minutes}
\newcommand{\UnitNumberName}{Unit 2: Exponential and Logarithmic Functions \quad}
\newcommand{\LessonNumberName}{Lesson 2.9: Logarithmic Expressions}

\begin{document}
\begin{center}
    {\Large\bfseries \CourseName: \SchoolYear \\
    {\small\bfseries \UnitNumberName \LessonNumberName}}
\end{center}

\begin{tcolorbox}[colback=sky, colframe=navy, boxrule=0.9pt, arc=2mm,
    left=2mm, right=2mm, top=1.5mm, bottom=1.5mm]
    \textbf{Primary Objective:} Students define and evaluate logarithms using the inverse
    relationship with exponentiation, convert between exponential and logarithmic form, evaluate
    common and natural logarithms with and without technology, and interpret logarithmic scales.
    \hfill\textit{\small Skill: AP 1.B}
\end{tcolorbox}

\renewcommand{\arraystretch}{1.6}
\begin{skillbox}[Priority Ideas \& Skills]{goldbox}
\begin{tabularx}{\linewidth}{@{}X X@{}}
\textbf{AP Skills} &
\textbf{Key Understandings} \\
\midrule
\textbf{AP 1.B} --- Express functions in analytically equivalent forms. &
\begin{itemize}[leftmargin=1.2em,itemsep=2pt,topsep=0pt]
  \item $\log_b c = a \Leftrightarrow b^a = c$ (where $b > 0$, $b \ne 1$). (EK 2.9.A.1)
  \item Some logarithm values are exact; others require technology. (EK 2.9.A.2)
  \item On a logarithmic scale, each unit is a multiplicative change by the base. (EK 2.9.A.3)
\end{itemize} \\
\end{tabularx}
\end{skillbox}

\begin{skillbox}[{Vocabulary, Concepts, \& Theorems}]{greenbox}
\renewcommand{\arraystretch}{1.4}
\begin{tabularx}{\linewidth}{@{} >{\bfseries\color{navy}}p{0.27\linewidth} X @{}}
logarithm            & The exponent to which a base must be raised to produce a given value:
                       $\log_b c = a$ means $b^a = c$. \\
common logarithm     & A logarithm with base 10; written $\log x$ when no base is shown. \\
natural logarithm    & A logarithm with base $e \approx 2.718$; written $\ln x$. \\
exponential form     & An equation written as $b^a = c$. \\
logarithmic form     & An equation written as $\log_b c = a$. \\
logarithmic scale    & A scale where each unit represents a multiplicative change by the base;
                       equal spacing corresponds to equal ratios, not equal differences. \\
\end{tabularx}
\end{skillbox}

\begin{skillbox}[Activate Prior Knowledge \& Spiral Review]{sky}
    Spiral review (warm-up) connects to Lesson 2.8 (inverse functions): students recall that
    an inverse ``undoes'' the original function. They identify the exponent $x$ in $2^x = 32$
    and $10^x = 1000$ by inspection, then complete a sentence frame describing what a logarithm
    asks. This primes the conceptual link: $\log_b c$ is the exponent the inverse of $b^x$ produces.
\end{skillbox}

\begin{skillbox}[Hook]{sky}
Write on the board: \textbf{``An earthquake of magnitude 6 is how much stronger than magnitude 5?''}
Take a vote (most students guess ``twice as strong'' or ``slightly stronger'').
Reveal: it is \textbf{10 times} stronger. The Richter scale is a \textit{logarithmic scale} ---
each whole-number step means a 10$\times$ increase in ground motion.

Ask: ``If the scale were linear, what would magnitude 10 mean compared to magnitude 1?''
(9 times stronger.) ``What does a logarithmic scale say?'' ($10^9 = $ one billion times stronger.)
This motivates why logarithms matter: they compress enormously wide ranges into human-readable numbers.
\end{skillbox}

\begin{skillbox}[Lesson]{sky}
\begin{multicols}{2}
\textbf{Part 1 (12 min): Definition and conversion.}

Logarithm definition:
\[
  \log_b c = a \quad\Leftrightarrow\quad b^a = c
\]
where $b > 0$, $b \ne 1$, $c > 0$.

Walk through conversion table (exponential $\leftrightarrow$ log form): at least five examples,
including $b^0 = 1$ and $b^1 = b$, giving $\log_b 1 = 0$ and $\log_b b = 1$.

Key question: ``What does $\log_2 8$ ask?'' --- ``What power of 2 gives 8?'' (Answer: 3.)

\columnbreak

\textbf{Part 2 (12 min): Common and natural log; technology.}

Common log: $\log x = \log_{10} x$ (base 10 when omitted).
Natural log: $\ln x = \log_e x$, where $e \approx 2.718$.

Exact values: $\log_{10} 1000 = 3$; $\ln e = 1$; $\ln 1 = 0$.

Technology estimate: $\log_2 7 \approx 2.807$ (between 2 and 3, since $2^2 = 4 < 7 < 8 = 2^3$).

Bracket without calculator: $3^2 = 9 < 20 < 27 = 3^3$, so $\log_3 20$ is between 2 and 3.

\textbf{Part 3 (10 min): Logarithmic scale.}

Number line with positions $10^0, 10^1, 10^2, 10^3$. Each unit = multiply by 10.

Halfway between $10^1$ and $10^2$ on a log scale? $10^{1.5} \approx 31.6$ --- not 55
(the arithmetic midpoint). Stress: equal spacing on a log scale means equal ratios.
\end{multicols}
\end{skillbox}

\begin{skillbox}[Active Monitoring]{redbox}
\begin{itemize}[leftmargin=1.2em,itemsep=2pt,topsep=0pt]
  \item After Part 1: cold-call ``What does $\log_5 25$ equal? How do you know?''
        Expect: 2, because $5^2 = 25$. Watch for students who confuse base with result.
  \item During Part 1: watch for students who write $b^a = c$ as $c^a = b$ (transposing base
        and result). Reinforce: the base of the log is the base of the power.
  \item After Part 2: confirm students know $\log$ (no base) means base 10, and $\ln$ means base $e$.
  \item During Part 3: ask ``Is 50 closer to 10 or to 100 on a log scale?''
        ($\log_{10} 50 \approx 1.70$, so it is closer to 100 than to 10.)
\end{itemize}
\end{skillbox}

\begin{skillbox}[Group Work \& Differentiation]{redbox}
\begin{itemize}[leftmargin=1.2em,itemsep=2pt,topsep=0pt]
  \item \textbf{Tier R (Remediate):} Conversion drill --- exponential to log form and back.
        Evaluate exact logs where the answer is a small whole number.
  \item \textbf{Tier A (Apply):} Mix of exact evaluations, technology estimates, bounding
        without a calculator, and interpreting the pH formula $\text{pH} = -\log_{10}[\text{H}^+]$.
  \item \textbf{Tier E (Extend):} Prove $\log_b b^n = n$ and $b^{\log_b n} = n$ from the
        definition; solve $\log_4 x = \tfrac{3}{2}$ by converting; connect fractional exponents
        to fractional logarithms.
\end{itemize}
\end{skillbox}

\begin{skillbox}[Individual Work \& Assessment]{redbox}
Exit ticket (3 questions, 1 page): (1) convert $3^5 = 243$ to log form and evaluate;
(2) evaluate $\log_{10} 10{,}000$ and $\ln e^3$ without a calculator;
(3) bound $\log_2 50$ between two consecutive integers and explain.

Sort by performance: students who miss (1) need more conversion practice (EK 2.9.A.1);
those who miss (2) may not have internalized $\log_b b^n = n$; those who miss (3) need
number-sense work on bracketing logarithms (EK 2.9.A.2).
\end{skillbox}

\begin{skillbox}[Reinforcement \& Extension]{goldbox}
Homework (6 problems): conversion practice (6 items); evaluate without a calculator (5 items);
technology estimates to 3 decimal places; bounding between integers without a calculator;
real-world decibel formula; AP-style MC on the definition.

Preview of Lesson 2.10: logarithmic functions as inverses of exponential functions ---
connecting the algebraic definition $\log_b c = a \Leftrightarrow b^a = c$ to the behavior
of $f(x) = \log_b x$ as a function: its domain, range, intercepts, and how its graph relates
to the graph of $g(x) = b^x$ by reflection across $y = x$.
\end{skillbox}

\end{document}
